% Partie : set-theory
% Chapitre : ordinals
% Section : wo

\documentclass[../../../include/open-logic-section]{subfiles}

\begin{document}

\olfileid[fr]{sth}{ordinals}{wo}
\olsection{Bons ordres}

La notion fondamentale est la suivante :

\begin{defn}
La relation $<$ \emph{définit un bon ordre} sur $A$ si et seulement si
elle satisfait aux deux conditions suivantes :
\begin{enumerate}
	\item $<$ est connexe, c'est-à-dire que, pour tous $a, b \in A$,
	soit $a < b$, soit $a = b$, soit $b < a$ ;
	\item tout sous-ensemble non vide de $A$ possède un !!{element}
	minimal pour $<$, c'est-à-dire que, si $\emptyset \neq X \subseteq A$,
	alors $(\exists m \in X)(\forall z \in X)z \nless m$.
\end{enumerate}
\end{defn}

On voit aisément que les trois exemples considérés à la section
précédente sont bien des relations de bon ordre.

\begin{prob}
La \olref[sth][ordinals][idea]{sec} présente trois exemples d'ordres
sur les nombres naturels. Vérifiez que chacun est un bon ordre.
\end{prob}

Voici quelques observations élémentaires, mais extrêmement importantes,
sur les bons ordres.

\begin{prop}\ollabel{wo:strictorder}
Si $<$ définit un bon ordre sur $A$, alors tout sous-ensemble non vide
de $A$ possède un unique plus petit élément pour $<$, et $<$ est
irréflexive, asymétrique et transitive.
\end{prop}

\begin{proof}
Si $X$ est un sous-ensemble non vide de $A$, il possède un !!{element}
minimal $m$ pour $<$, c'est-à-dire que $(\forall z \in X)z \nless m$.
Comme $<$ est connexe, $(\forall z \in X)m \leq z$. Ainsi, $m$ est le
plus petit !!{element} de $X$ pour $<$.

Pour l'irréflexivité, fixons $a \in A$ : le plus petit !!{element} de
$\{a\}$ pour $<$ est $a$, donc $a \nless a$. Pour la transitivité,
si $a < b < c$, alors, puisque $\{a, b, c\}$ possède un plus petit
!!{element} pour $<$, on a $a < c$. L'asymétrie découle de
l'irréflexivité et de la transitivité.
\end{proof}

\begin{prop}\ollabel{propwoinduction}
Si $<$ définit un bon ordre sur $A$, alors, pour toute formule $\phi(x)$ :
\[
	\text{si }(\forall a \in A)((\forall b < a)\phi(b) \lif
		\phi(a))\text{, alors }(\forall a \in A)\phi(a).
\]
\end{prop}

\begin{proof}
Nous allons démontrer la contraposée. Supposons que
$\lnot(\forall a \in A)\phi(a)$, c'est-à-dire que
$X = \Setabs{x \in A}{\lnot\phi(x)} \neq \emptyset$. Alors $X$
possède un !!{element} minimal pour $<$, disons $a$. Nous avons donc
$(\forall b < a)\phi(b)$, mais $\lnot \phi(a)$.
\end{proof}\noindent Cette dernière propriété devrait rappeler le
principe de récurrence forte sur les nombres naturels : si
$(\forall n \in \omega)((\forall m < n)\phi(m) \lif \phi(n))$, alors
$(\forall n \in \omega)\phi(n)$. C'est cette propriété qui fait de
la notion de bon ordre une notion très \emph{robuste}.\footnote{Rappelons
que toutes les formules peuvent contenir des paramètres, sauf indication
contraire explicite.}

\end{document}
